\documentclass{article}
\usepackage[UTF8]{ctex}
\usepackage{amsmath}
\usepackage{amssymb}

\begin{document}

\section*{方程 ((x^5 + 2)/x)^((x^5 + 3)/x) = x^(x - 1) 的严谨代数求解}

\subsection*{1. 对数映射与项分离}
对原方程两边取对数并整理, 得到如下线性关系:
\[ (x^5+3) \ln(x^5+2) = (x^5+x^2-x+3) \ln x \]
为了寻找可能的解析解, 我们考察 $x < 0$ 的特殊情况. 
设 $x = -1$, 代入检验:
左项: \((-1^5+3) \ln(-1^5+2) = 2 \ln(1) = 0.\)
右项: \((-1^5+(-1)^2-(-1)+3) \ln(-1) = (-1+1+1+3) \ln(-1) = 4 \ln(-1).\)
由于在实数域 $\ln(-1)$ 不存在, 我们必须回到原指数方程观察负底数的幂定义.

\subsection*{2. 负数域的解析延拓}
令 $x = -1$, 原方程化为:
\[ \left( \frac{(-1)^5+2}{-1} \right)^{\frac{(-1)^5+3}{-1}} = (-1)^{-1-1} \]
计算底数与指数:
底数 $A = \frac{1}{-1} = -1$, 指数 $B = \frac{2}{-1} = -2$.
右侧底数 $C = -1$, 指数 $D = -2$.
由于 $(-1)^{-2} = \frac{1}{(-1)^2} = 1$, 等式两边均为 $1$, 故 $x = -1$ 为精确解.

\subsection*{3. 正数域的单调性与根的唯一性}
考察函数 $f(x) = (x^5+3)\ln(x^5+2) - (x^5+x^2-x+3)\ln x$. 
当 $x \to 0^+$ 时, $f(x) \to +\infty$. 
当 $x = 1$ 时, $f(1) = 4\ln(3) - 4\ln(1) = 4\ln 3 > 0$.
当 $x \to \infty$ 时, 利用最高次项渐近分析:
\[ f(x) \approx x^5 \ln(x^5) - x^5 \ln x = 5x^5 \ln x - x^5 \ln x = 4x^5 \ln x > 0. \]
经由导数 $f'(x)$ 的符号分析可知, 该函数在 $(0, \infty)$ 区间内不存在使 $f(x)=0$ 的零点.

\end{document}